\section{Evaluation Protocol}
\label{sec:protocol}

\subsection{Baseline Families}

CLAIMFORGE should be evaluated against methods from six families. The first family is image manipulation localization, using a unified harness such as IMDL-BenCo where possible \citep{ma2024imdldbenco}. This family includes dual-output or localization-first models such as TruFor, HiFi-IFDL, Mesorch, and SAFIRE \citep{guillaro2022trufor,guo2023hifi,zhu2024mesorch,kwon2024safire}. The second family is whole-image AI-generated image detection, including detectors such as DIRE and C2P-CLIP \citep{wang2023dire,tan2024c2pclip}. The third family is diffusion-inpainting-specific detection and localization, represented by GIM, OpenSDI-style models, and COCO-Inpaint-style protocols \citep{chen2024gim,wang2025opensdi,yan2025cocoinpaint}. The fourth family is multimodal or vision-language forensic reasoning: these models can be prompted to identify whether the image evidence is plausible and to describe suspicious regions. The fifth family is human and domain-expert review, which measures whether the images are deceptive in the claim setting. The sixth family is active provenance and watermarking, which is treated as an orthogonal defense rather than a passive detector baseline.

\subsection{Metrics}

Image-level detection reports AUC, AP, and fixed-threshold accuracy. AUC and AP are mandatory because fixed thresholds can understate or overstate separability when models are miscalibrated \citep{yang2026calibration}. Localization is scored on forged images using pixel F1 at a fixed threshold, best-threshold F1, IoU, MCC, and pixel AP. For methods that output both an image score and a mask, we also report a detect-then-localize composite:
\[
S_{\mathrm{joint}} = \mathrm{F1}_{\mathrm{image}} \times \mathrm{F1}_{\mathrm{pixel}}.
\]
This composite is not meant to replace the two leaderboards. It prevents a detector from ranking highly by solving only detection or only localization.

\subsection{Generalization Regimes}

The main benchmark claim should be tested under three regimes:
\begin{enumerate}
    \item \textbf{Zero-shot.} Evaluate off-the-shelf weights trained on external data.
    \item \textbf{In-domain fine-tuning.} Fine-tune on CLAIMFORGE train and test on matched held-out samples.
    \item \textbf{Held-out editor, held-out domain, and laundering.} Fine-tune on CLAIMFORGE train, then test on an unseen editor, unseen domain, and post-processing transformations.
\end{enumerate}
The third regime is the headline stress test because it matches the intended deployment gap. A method should be said to ``solve'' CLAIMFORGE only if it reaches both high image-level detection and usable localization under this regime. A concrete bar for the final paper is image AUC at least 0.90 and pixel F1 at least 0.50 on the held-out, laundered test set. This bar should be stated before running baselines.

\subsection{Shortcut Audits}

The benchmark should include explicit shortcut audits. A format-only probe should test whether file metadata, compression statistics, or image size can separate real from forged images. A region-prior probe should test whether insertion locations leak class information. A real-object paste-back control should test whether localizers are detecting paste artifacts rather than AI-generated texture. Finally, laundering curves should report how each method changes under JPEG recompression, resizing, blur/noise, screenshotting, and social-media round trips.
